% !TEX root = chapter-standalone.tex

\section{Partitions}
\label{sec:partitions}

\begin{marginfigure}
    \centering
    \includegraphics[width=\linewidth]{set_of_foods}
    \caption{A set of foods.}
    \label{fig:set_of_foods}
\end{marginfigure}
Let's start with a simple toy example to illustrate the idea of grouping together elements of a set into a partition. Consider a set that consists of 10 different foods in total, and is made up of three bread-like foods, two desserts, four fruits, and one cheese. Now suppose that we want to tell someone only the main \emph{kinds} of food that are in our set. Then we don't need to include any information about how many fruits, desserts, etc. there are, we only want to communicate the fact that there are the following four kinds of food: bread-like, desserts, fruit, and cheese. 

\begin{marginfigure}
    \centering
    \includegraphics[width=\linewidth]{partition_kinds_of_food}
    \caption{A partition of a set of foods.}
    \label{fig:partition_kinds_of_foods}
\end{marginfigure}
On a mathematical level, one way to go from the ten-element set of \emph{foods} to a four-element set of \emph{kinds of foods} is to partition the ten-element set of foods into disjoint subsets, where each subset contains only the foods of a given kind. Then we can think of these four subsets as \emph{representing} or \emph{signifying} the four kinds of food in question (and we can simply ignore how large each of these subsets is, because that is not relevant to the question of how many \emph{kinds} there are). In other words, the four subsets in the partition can play the role of ``names'' for the four kinds of food. 

%Furthermore, the four subsets representing our ``kinds'' do form a set, namely this one: 
%
%Thus our construction produces a four-element set from our original ten-element set. 

Let's start to make the ideas in the above example more precise by defining partitions. Given any set $\setA$, a partition of $\setA$ is a collection of subsets of $\setA$ which, on the one hand, are pairwise disjoint (they don't intersect) and on the other hand the union of all the subsets in the collection is equal to $\setA$. In other words, the collection of subsets ``covers'' all of $\setA$. 

\begin{ctdefinition}[Partition]
    \label{def:partition}
    A \maindef{partition} of a set~\setA is
    
    \constit
    \begin{enumerate}
    \item A set~$\makeset{\setA_\setIel}_{\setIel\setin \setI}$ of subsets~$\setA_\setIel\setsubseteq \setA$
    \end{enumerate} 
    \condit
    \begin{enumerate}
        \item If $\setIel\neq \setJel$,  then $\setA_\setIel\setintersection \setA_\setJel=\Emptyset$;
        \item $\bigsetunion_{\setIel\setin \setI}\setA_\setIel=\setA$.
    \end{enumerate}
\end{ctdefinition}



\todotext{The emojis for this example need to be set up in symbols.tex...}
\begin{example}
If $\setA$ denotes the set of foods depicted in \cref{fig:set_of_foods}, then an example of a partition of this set is
\begin{itemize}
\item $\setA_1 = \makeset{A}$
\item $\setA_2 = \makeset{B}$
\item $\setA_3 = \makeset{C}$
\item $\setA_4 = \makeset{D}$
\end{itemize}
\end{example}

\begin{marginfigure}
    \centering
    \includegraphics[width=\linewidth]{partition_another_kinds_of_food}
    \caption{Another partition of a set of foods.}
    \label{fig:partition_another_kinds_of_foods}
\end{marginfigure}
Note that there are of course, in general, many different possible partitions of a given set. For the set of foods in \cref{fig:set_of_foods}, we could also consider the partition shown in \cref{fig:partition_another_kinds_of_foods} This partition might correspond to the idea of grouping the foods into the three categories ``baked goods'', ``healthy snacks'', and ``candy''.




